% Live talk on the SHiP veto study, for PictureThis (widgets computed by veto_talk.py). Build and present: ./run.sh
% In any other PDF viewer the widgets are grey placeholders.
\documentclass[aspectratio=169,10pt]{beamer}
\usepackage{talktheme}
\usepackage{picturethis}
\hypersetup{hidelinks}
\usetikzlibrary{decorations.pathmorphing}

\title[A veto for SHiP, live]{A veto for SHiP}
\subtitle{Muon background, design space and reconstruction, live}
\author{Shah Rukh Qasim}
\institute{University of Zurich · with G.~Frisella}
\date{October 2026}

\tikzset{
  box/.style={draw=line, fill=white, rounded corners=3pt, inner sep=4pt, align=center, font=\footnotesize},
  hbox/.style={box, fill=wash},
  abox/.style={box, draw=accent, line width=0.8pt},
  arr/.style={-{Stealth[length=2.2mm]}, line width=0.8pt, color=muted},
  aarr/.style={-{Stealth[length=2.2mm]}, line width=0.9pt, color=accent},
}
\newcommand{\sld}[2][\linewidth]{\livewidget{#2}{#1}{0.62cm}}
\newcommand{\legendline}[1]{{\color{muted}\lightfont\scriptsize #1}}
\newcommand{\sw}[1]{\textcolor[HTML]{#1}{\rule{1.1ex}{1.1ex}}}

\begin{document}

\begin{frame}[plain,noframenumbering]
  \titlepage
\end{frame}

\begin{frame}{How to read this talk}
  \begin{columns}[c]
    \begin{column}{0.56\textwidth}\small
      Every plot with a slider is computed while you watch, by the code of the study (\texttt{veto/lib/veto}):
      \begin{itemize}
        \item muons after the shield from our own production (the enriched input), and the older samples
        \item interactions of muons and neutrinos in everything after the shield (fast-veto-sim)
        \item digitization, random hits, the selection and the veto
        \item designs scored on every background source, by hand and by an LLM-driven evolution
      \end{itemize}
      \vspace{1mm}
      Slow widgets show a spinner and fill in. 3D views rotate and zoom.\\[1mm]
      \takeaway[\footnotesize]{The spill is \key{1 s} of $5\cdot10^{13}$ POT and the weights are muons per spill, so muons per spill and Hz are the same number. 15 years $= 1.2\cdot10^{7}$ spills.}
    \end{column}
    \begin{column}{0.4\textwidth}\centering
      \begin{tikzpicture}[node distance=1.6mm, every node/.append style={font=\scriptsize}]
        \node[hbox, text width=3.9cm] (data) {\semibold muons after the shield\\\faint{our production, enriched input}};
        \node[hbox, below=of data, text width=3.9cm] (int) {\semibold interactions\\\faint{Geant4 + Pythia\,6 ($\mu$), GENIE ($\nu$)}};
        \node[box, below=of int, text width=3.9cm] (des) {\semibold a veto design\\\faint{UBT + fins, JSON}};
        \node[abox, below=of des, text width=3.9cm] (reco) {\semibold digitize, select, veto\\\faint{backward search with a vertex cone}};
        \node[box, below=of reco, text width=3.9cm] (score) {\semibold score\\\faint{every background, signal, cost, rates}};
        \node[box, below=of score, text width=3.9cm] (rl) {\semibold optimise\\\faint{an LLM evolves the design}};
        \draw[arr] (data) -- (int); \draw[arr] (int) -- (des); \draw[arr] (des) -- (reco); \draw[aarr] (reco) -- (score); \draw[arr] (score) -- (rl);
      \end{tikzpicture}
    \end{column}
  \end{columns}
\end{frame}

% ======================================================================== PART 1
\blurb{Muons from the target reach the detector. Two of them can fake a decay in the vessel, and one muon or neutrino interacting after the shield makes a real vertex. The veto has to see them coming in, without killing the signal.}
\section{The problem}

\begin{frame}{SHiP and its muon background}
  \begin{columns}[c]
    \begin{column}{0.68\textwidth}
      \livewidget{det3d}{\linewidth}{0.6\textheight}\\[1mm]
      \sld[0.6\linewidth]{n_show}
    \end{column}
    \begin{column}{0.3\textwidth}\footnotesize
      \key{Target} at $z=0$, \key{muon shield} (TRY6) to 30\,m.\\[1mm]
      \key{Decay vessel} $z = 30.6$--$80.6$\,m, $1\times2.7$\,m entrance widening to $4\times6$\,m.\\[1mm]
      \key{UBT} on the entrance window, \key{SST1} at 82\,m, spectrometer magnet and timing detector behind.\\[1mm]
      Cavern wall at $x=-4.56$\,m, floor at $y=-3.36$\,m.\\[2mm]
      \legendline{Muons per type (slider): \sw{9400D3}\,bounced into SST1 \sw{D62728}\,SST1 $p>1$, timing \sw{FF9896}\,SST1 other \sw{1F77B4}\,through an SBT wall \sw{FF7F0E}\,through the UBT window \sw{E377C2}\,bounced in the rock \sw{7F7F7F}\,straight past \sw{8C564B}\,stop in the rock \sw{2B2B2B}\,stop in the shield.}
    \end{column}
  \end{columns}
\end{frame}

\begin{frame}{What the veto has to do}
  \begin{columns}[c]
    \begin{column}{0.5\textwidth}\centering
      \begin{tikzpicture}[x=0.95cm,y=0.9cm]
        \fill[wash] (0,-1) -- (4.6,-1.9) -- (4.6,1.9) -- (0,1) -- cycle;
        \draw[accent, line width=0.8pt] (0,-1) -- (4.6,-1.9) (0,1) -- (4.6,1.9);
        \node[font=\scriptsize\lightfont, color=accent] at (3.6,-2.05) {decay vessel};
        \draw[ora, line width=2pt] (0,-1) -- (0,1); \node[font=\scriptsize, color=ora, below] at (0,-1) {UBT};
        \draw[grn, line width=2pt] (1.0,-1.2) -- (1.0,-2.4); \node[font=\scriptsize, color=grn, left] at (1.0,-2.3) {fin};
        \draw[red, line width=1pt] (-0.8,-2.2) -- (1.6,-1.35) -- (5.2,0.6);
        \draw[red, line width=1pt] (-0.8,0.3) -- (5.2,-0.9);
        \fill[red] (3.05,-0.26) circle (0.07); \node[font=\scriptsize, color=red, above] at (3.05,-0.2) {fake vertex};
        \fill[white, draw=grn, line width=1pt] (1.0,-1.89) circle (0.07);
        \fill[white, draw=ora, line width=1pt] (0,0.14) circle (0.07);
        \draw[ink, line width=1pt, decorate, decoration={snake, amplitude=1pt, segment length=4pt}] (-0.8,0.12) -- (2.6,0.25);
        \fill[ink] (2.6,0.25) circle (0.06);
        \draw[ink, line width=1pt] (2.6,0.25) -- (5.2,1.3) (2.6,0.25) -- (5.2,-0.2);
        \node[font=\scriptsize, right, fill=white, inner sep=1pt] at (-0.45,1.55) {signal (wavy): neutral from the target, no hits};
        \draw[muted] (5.2,-2.0) -- (5.2,2.0); \node[font=\scriptsize, color=muted, right] at (5.2,1.9) {SST1};
      \end{tikzpicture}
    \end{column}
    \begin{column}{0.47\textwidth}\small
      Two muons from the same spill, within $\pm200$\,ps, can form a vertex inside the vessel that points back to the target: a \key{fake signal}.\\[2mm]
      Each came in from outside: through the UBT window, the side walls, or after bouncing off the cavern. A hit where its track came from \key{explains} it.\\[2mm]
      \begin{tabular}{@{}ll@{}}
        muons per spill (simulated) & $1.8\cdot10^{10}$\\
        crossing $z = 30.6$\,m & $3.5\cdot10^{9}$\\
        in the SST1 acceptance & $3\cdot10^{4}$\\
        fake-vertex pairs, 15 years & \num{36} before any veto\\
      \end{tabular}\\[2mm]
      \takeaway[\footnotesize]{Veto every fake vertex, but do not veto the signal with \key{random hits} of the $10^{9}$--$10^{10}$ other muons.}
    \end{column}
  \end{columns}
\end{frame}

\begin{frame}{Three kinds of background}
  \footnotesize
  \begin{tabular}{@{}p{0.2\linewidth}p{0.24\linewidth}p{0.24\linewidth}p{0.24\linewidth}@{}}
    \toprule
    & \key{combinatorial} & \key{muon interactions} & \key{neutrino interactions}\\
    \midrule
    origin & two unrelated muons within $\pm200$\,ps & one muon interacts in the helium, the walls, the air or a veto plane & one neutrino interacts in the same materials\\[1mm]
    vertex & fake (a coincidence) & real & real\\[1mm]
    can the veto see it? & yes: each muon came in from outside & yes in principle: the parent crossed the UBT or a fin, but it was bent at the vertex & no: nothing charged comes in\\[1mm]
    handled by & the veto & the veto with a \key{vertex cone} & the \key{selection} only\\
    \bottomrule
  \end{tabular}\\[3mm]
  \takeaway[\footnotesize]{The first studies only had the first column. The other two are larger by orders of magnitude before any cut; they changed the selection, the veto and the objective.}
\end{frame}

% ======================================================================== PART 2
\blurb{Before any statistics: what a single muon does after the shield, where it crosses the detector, how it bounces, scatters and interacts.}
\section{One particle at a time}

\begin{frame}{One muon at a time}
  \begin{columns}[T]
    \begin{column}{0.27\textwidth}
      \livewidget{mu_type}{\linewidth}{6.2cm}
    \end{column}
    \begin{column}{0.71\textwidth}
      \livewidget{mu3d}{\linewidth}{0.6\textheight}\\[1mm]
      \sld{mu_index}\\[0.3mm]
      \legendline{15 typical muons per type (full run, drawn by weight); straight lines between the stored states. \textbullet\ $z = 30.6$\,m, $\blacksquare$ $z = 82$\,m or stop, \texttimes\ bounce.}
    \end{column}
  \end{columns}
\end{frame}

\begin{frame}{Bouncing off the cavern}
  \livewidget{bounce_view}{0.6\linewidth}{0.56\textheight}\hfill\livewidget{bounce_map}{0.38\linewidth}{0.56\textheight}\\[1mm]
  \sld[0.4\linewidth]{bounce_ev}\\[0.5mm]
  \legendline{Grey dashed: straight line from the 30.6\,m state; orange dotted: backward line from the 82\,m state. They meet within 4\,cm (median) at the bounce point, on the $-x$ wall or the floor. Right: all reconstructed bounce points trace the cavern.}
\end{frame}

\begin{frame}{One muon at a time: what it crosses}
  \livewidget{sim_mu_one3d}{\linewidth}{0.58\textheight}\\[1mm]
  \begin{columns}[c]
    \begin{column}{0.44\textwidth}\sld{sim_mu_pick}\end{column}
    \begin{column}{0.52\textwidth}\livewidget{sim_mu_design}{\linewidth}{0.62cm}\end{column}
  \end{columns}
  \legendline{Muons that do interact somewhere. Blue: helium chord · red: PVC (exact chord) · grey: air · dots: planes. $P = n_\mathrm{area}\,\sigma(E)$, the chance to interact there.}
\end{frame}

\begin{frame}{Muon--nucleus interactions}
  \begin{columns}[c]
    \begin{column}{0.44\textwidth}\centering
      \begin{tikzpicture}[x=1cm,y=1cm]
        \draw[red, line width=1.2pt, -{Stealth}] (0,2.2) -- (1.6,1.4) node[midway, above, font=\scriptsize] {$\mu\,(E)$};
        \draw[red, line width=1.2pt, -{Stealth}] (1.6,1.4) -- (3.6,2.3) node[midway, above, font=\scriptsize] {$\mu'\,(E-\nu)$};
        \draw[accent, line width=1pt, decorate, decoration={snake, amplitude=2pt, segment length=6pt}] (1.6,1.4) -- (1.9,0.2);
        \node[font=\scriptsize, color=accent] at (2.45,0.85) {$\gamma^*\ (\nu, Q^2)$};
        \draw[ink, line width=1.6pt] (0,-0.2) -- (1.9,0.0); \node[font=\scriptsize, left] at (0,-0.2) {N};
        \fill[ora!60] (1.9,0.0) circle (0.28);
        \foreach \a in {-25,-10,5,18,32}{\draw[ora, line width=0.8pt, -{Stealth}] (2.15,0.0) -- ++(\a:1.5);}
        \node[font=\scriptsize, color=ora] at (3.85,0.1) {hadrons, $W$};
      \end{tikzpicture}\\[2mm]
      {\footnotesize $W^2 = M^2 + 2M\nu - Q^2$}
    \end{column}
    \begin{column}{0.53\textwidth}\small
      The muon exchanges a virtual photon with a nucleus: the muon goes on, \key{bent by a few mrad}, and hadrons
      come out (pions, kaons, protons, K$^0_S$, $\Lambda$ that decay later).\\[1.5mm]
      \key{Geant4} (muonNuclear): the soft part, resonances and nuclear break-up. \key{Pythia\,6}: the hard part, $W > 2$\,GeV.
      \key{crossover} (our choice): Geant4 below an energy transfer $\nu_x(E)$, Pythia above.\\[1.5mm]
      \begin{tabular}{@{}lrrr@{}}
        $E_\mu$ [GeV] & 10 & 100 & 300\\\midrule
        $\sigma$ per nucleon [$\mu$b] & 8.7 & 19.8 & 26.2\\
        of which Pythia & 1\,\% & 11\,\% & 15\,\%\\
      \end{tabular}\\[1mm]
      \aside{Library: 50 energies 0.5--500\,GeV, 20\,k events per energy and material.}
    \end{column}
  \end{columns}
\end{frame}

\begin{frame}{One interaction at a time}
  \livewidget{sim_int3d}{\linewidth}{0.56\textheight}\\[1mm]
  \begin{columns}[c]
    \begin{column}{0.3\textwidth}\sld{sim_i_E}\end{column}
    \begin{column}{0.3\textwidth}\sld{sim_i_ev}\end{column}
    \begin{column}{0.36\textwidth}\sld{sim_i_zoom}\end{column}
  \end{columns}
  \livewidget{sim_i_gen}{0.45\linewidth}{0.62cm}\quad
  \legendline{Library events: beam along $z$; length $\propto\log p$; thick: the scattered muon; dashed: neutral; $\diamond$: decay point. E: 3, 10, 30, 100, 300\,GeV.}
\end{frame}

\begin{frame}{One applied interaction at a time}
  \livewidget{sim_ai3d}{\linewidth}{0.58\textheight}\\[1mm]
  \begin{columns}[c]
    \begin{column}{0.32\textwidth}\sld{sim_v_ev}\end{column}
    \begin{column}{0.66\textwidth}\livewidget{sim_v_tgt}{0.72\linewidth}{0.62cm}\ \livewidget{sim_v_kind}{0.26\linewidth}{0.62cm}\end{column}
  \end{columns}
  \legendline{Black: the parent muon (A, K, B); dotted grey: the parent neutrino. Products straight to $z = 82$\,m. "→ SST1": charged, $p > 1$\,GeV, in $|x|<2$, $|y|<3$\,m.}
\end{frame}

\begin{frame}{Decays: from where, into what}
  \livewidget{sim_decay3d}{\linewidth}{0.58\textheight}\\[1mm]
  \begin{columns}[c]
    \begin{column}{0.55\textwidth}\livewidget{sim_dk_part}{\linewidth}{0.62cm}\end{column}
    \begin{column}{0.42\textwidth}\sld{sim_dk_p}\end{column}
  \end{columns}
  \legendline{K$^0_S$, hyperons, $\pi^0$, $\eta$ are decayed in the library (main channels, decay point from $c\tau$): a K$^0_S$ made in the wall can decay inside the vessel into $\pi^+\pi^-$, a real two-track vertex.}
\end{frame}

\begin{frame}{Multiple scattering: one muon at a time}
  \livewidget{sim_msc_traj}{\linewidth}{0.62\textheight}\\[1mm]
  \sld[0.5\linewidth]{sim_ms_traj_p}\quad\legendline{60 muons on the axis; red: a hard single scatter. 1.5, 3, 10, 30, 100\,GeV.}
\end{frame}

% ======================================================================== PART 3
\blurb{Our own production of the muons after the shield, an input enriched at high energy, and the forced interactions of muons and neutrinos on top. Simulated once; every design is then evaluated in seconds.}
\section{Data and simulation}

\begin{frame}{Two simulations of the same muons}
  \footnotesize
  \begin{tabular}{@{}>{\raggedright}p{0.14\textwidth}>{\raggedright}p{0.38\textwidth}>{\raggedright\arraybackslash}p{0.4\textwidth}@{}}
    \toprule
    & \key{20\,001 old runs} & \key{the full run} (\texttt{Full\_dataset\_z\_3060\_8200.pkl})\\
    \midrule
    muons & only those ending in the SST1 acceptance, $\approx1\,300$ per run, 26.5\,M rows & all 506\,M simulated, one row each, 32\,GB\\
    states & start, $z = 30.6$\,m, $z = 82$\,m & $z = 30.6$\,m and $z = 82$\,m, or where the muon stopped\\
    good for & pair background: $\approx2\cdot10^{5}$ fake-vertex candidates & everything that misses the spectrometer: UBT and fin rates, veto-only muons\\
    read with & a column store, 78\,s to build & memory-mapped in place, a pass in $\approx1$\,min\\
    \bottomrule
  \end{tabular}\\[2mm]
  \begin{itemize}\footnotesize
    \item The full run is a \key{rerun of seed 42}: same input muons, same states at 30.6\,m; old file 42 holds its SST1 muons. Combined, it is not a 20\,002nd run.
    \item \texttt{evt\_n} = input-muon index in both; weights belong to the input muon (identical in both).
    \item Muons that bounce off the cavern differ between the two: bouncing is random from run to run.
  \end{itemize}
\end{frame}

\begin{frame}{Where the muons go}
  \begin{columns}[c]
    \begin{column}{0.62\textwidth}
      \livewidget{where_bar}{\linewidth}{0.7\textheight}
    \end{column}
    \begin{column}{0.36\textwidth}\footnotesize
      70\,\% stop in the shield; 26\,M range out in the rock behind the cavern wall.\\[1.5mm]
      \key{125\,M} pass the vessel to $z = 82$\,m outside the spectrometer: harmless themselves, but they light up any detector placed in their way.\\[1.5mm]
      Only $\approx2\,000$ per run matter directly: they enter the vessel or reach SST1.\\[1.5mm]
      \takeaway[\footnotesize]{The veto sees $10^{9}$--$10^{10}$ muons per spill to catch $\approx10^{4}$. Accidental hits are the real enemy.}
    \end{column}
  \end{columns}
\end{frame}

\begin{frame}{Where the muons go, in Hz}
  \begin{columns}[c]
    \begin{column}{0.62\textwidth}
      \livewidget{where_hz}{\linewidth}{0.7\textheight}
    \end{column}
    \begin{column}{0.36\textwidth}\footnotesize
      The same muons as rates: muons per spill of $5\cdot10^{13}$ POT over the 1\,s spill.\\[1.5mm]
      \key{$\approx$3.5\,GHz} cross $z = 30.6$\,m and 2.7\,GHz pass the vessel to $z = 82$\,m.\\[1.5mm]
      $\approx$29\,kHz reach the SST1 acceptance (the three SST1 bars together), $\approx$27\,kHz enter the vessel without reaching it.\\[1.5mm]
      \takeaway[\footnotesize]{A veto detector sees GHz in total; the muons it must catch arrive at tens of kHz.}
    \end{column}
  \end{columns}
\end{frame}

\begin{frame}{Our own muon production}
  \begin{columns}[c]
    \begin{column}{0.52\textwidth}\footnotesize
      \texttt{bin/generate\_muon\_study.py}: G.~Frisella's shield simulation (TRY6, GPU propagation) imported unchanged; every muon taken to $z = 30.6$\,m and 82\,m and classified on the GPU (torch, identical to the CPU version, $2\times151$\,M muons checked).\\[1.5mm]
      \key{One seed = one spill}: the same input muons, new random numbers. $\approx95$\,s per seed on one GPU.\\[1.5mm]
      \key{Every category is kept} up to 15\,M rows, also the muons that stop in the shield. Storage is \emph{fill until the budget}: type $t$ (rate $R_t$ per seed) is kept by the first
      \[ n_t = \max\big(\lceil 15\,\text{M}/R_t\rceil,\ 20\big) \]
      seeds, each a random fraction re-weighted to the seed's total. SST1 and veto-only muons: all of them, in every seed.
    \end{column}
    \begin{column}{0.45\textwidth}\scriptsize
      \begin{tabular}{@{}lrr@{}}
        \toprule
        type & per seed & seeds kept\\
        \midrule
        0 SST1, bounced & 299 & all\\
        1 SST1, selected & 301 & all\\
        2 SST1, other & 723 & all\\
        3 veto only (walls, window) & 975 & all\\
        4 bounced in the rock & 575\,k & 27\\
        5 straight past the vessel & 125\,M & 20\\
        6 stops at 30.6--82\,m & 26\,M & 20\\
        7 stops in the shield & 355\,M & 20\\
        \bottomrule
      \end{tabular}\\[1.5mm]
      \footnotesize Guglielmo's study files had 1\,500 picks per category, overlapping categories and a window cut too narrow for the UBT: the reason for our own production.\\[1mm]
      \aside{20\,000 seeds on farm-gpu3, 4 GPUs (PR \#7).}
    \end{column}
  \end{columns}
\end{frame}

\begin{frame}{Three datasets}
  \footnotesize
  \begin{tabular}{@{}l>{\raggedright}p{0.4\textwidth}rr>{\raggedright\arraybackslash}p{0.2\textwidth}@{}}
    \toprule
    & built from & rows & size & used for\\
    \midrule
    \key{alpha} (14 Oct) & 20\,001 old runs (SST1 types) + the full run (everything else) & 57\,M & 5.8\,GB & first results, first evolution\\
    \key{gamma} (17 Oct) & SST1 types from alpha; types 3--7 from 1\,715 seeds of our production & 88\,M & 8.9\,GB & interactions, cone\\
    \key{delta} (19 Oct) & everything from 484 seeds of the \key{enriched} production (input with the Pythia\,6 tail) & 111\,M & 11\,GB & \key{the end-to-end results}\\
    \bottomrule
  \end{tabular}\\[3mm]
  \begin{columns}[T]
    \begin{column}{0.5\textwidth}\footnotesize
      Same columns in all three: states at 30.6 and 82\,m, start state, bounce point, walls crossed, type, seed, \texttt{evt\_n} (the input muon). \texttt{weight} is always muons per spill: a seed's rows of a type are re-weighted to its total and divided by the number of seeds that kept the type.
    \end{column}
    \begin{column}{0.46\textwidth}\footnotesize
      Rows are not independent: every seed re-simulates the same input muons. What counts is the number of \key{independent input muons}, $n_\text{eff} = (\sum w)^2/\sum w^2$ with the weights summed per \texttt{evt\_n} first.\\[1mm]
      \takeaway[\footnotesize]{More seeds of the same input do not help where a few heavy input muons dominate.}
    \end{column}
  \end{columns}
\end{frame}

\begin{frame}{The enriched input: a Pythia\,6 tail}
  \begin{columns}[c]
    \begin{column}{0.6\textwidth}
      \livewidget{ds_neff}{\linewidth}{0.66\textheight}
    \end{column}
    \begin{column}{0.38\textwidth}\footnotesize
      In the 506\,M input muons, everything above 150\,GeV from minimum bias is a few rows of weight 768.75 each: in gamma \key{10 input muons carry 32\,\%} of the selected SST1 muons ($n_\text{eff} = 55$).\\[1.5mm]
      \texttt{fast-veto-sim} replaces them by \key{Pythia\,6.4.28} muons (first collision, Geant4 cascade, $\pi$/K forced to decay, Mo+W target), normalised to the same flux above 150\,GeV: 601\,M rows.\\[1.5mm]
      Selected SST1: $n_\text{eff}$ 55 $\to$ \key{18\,200}. The harder spectrum also gives \key{40--55\,\% more} SST1 muons.\\[1.5mm]
      \takeaway[\footnotesize]{Veto-only muons (walls, window) are unchanged at $\approx12$\,k: not the high-energy tail.}
    \end{column}
  \end{columns}
\end{frame}

\begin{frame}{Enough statistics everywhere?}
  \begin{columns}[c]
    \begin{column}{0.6\textwidth}
      \livewidget{probe3d}{\linewidth}{0.6\textheight}\\[1mm]
      \livewidget{probe_ds}{0.7\linewidth}{0.62cm}
    \end{column}
    \begin{column}{0.38\textwidth}
      \livewidget{probe_bands}{\linewidth}{0.36\textheight}\\[1mm]
      \footnotesize A probe, not a design: the UBT at full height, 42 fins every 2.5\,m and 5 floor planes tilted by 30$^\circ$; per 25\,cm cell the independent input muons crossing it.\\[1mm]
      \key{Thin}: the first metre from the vessel wall (median 11--14), the floor beyond 32\,m (8--13). Enrichment does not help there.\\[1mm]
      \aside{Red dotted: 10. Fins beyond 1\,m and the UBT: 140--430.}
    \end{column}
  \end{columns}
\end{frame}

\begin{frame}{From the shield to SST1: the chain}
  \begin{columns}[c]
    \begin{column}{0.54\textwidth}\centering
      \begin{tikzpicture}[node distance=1.6mm, every node/.append style={font=\scriptsize}]
        \node[hbox, text width=6.3cm] (mu) {\semibold muons after the shield\\\faint{the dataset (delta: 111\,M rows): A at 30.6\,m, bounce K, B at 82\,m}};
        \node[box, below=of mu, text width=6.3cm] (cross) {\semibold what each muon crosses\\\faint{helium chord, exact PVC chords, air, every plane: $P = n_\mathrm{area}\,\sigma(E)$}};
        \node[abox, below=of cross, text width=6.3cm] (force) {\semibold forced interactions\\\faint{draw $M$ crossings $\propto wP$, weight $S/M$; final state from a library (Geant4 + Pythia\,6; GENIE for $\nu$)}};
        \node[box, below=of force, text width=6.3cm] (tracks) {\semibold tracks to SST1\\\faint{straight lines, scattering with tails in walls, planes, helium, air}};
        \node[hbox, below=of tracks, text width=6.3cm] (veto) {\semibold our selection and veto, as for the signal};
        \draw[arr] (mu) -- (cross); \draw[arr] (cross) -- (force); \draw[aarr] (force) -- (tracks); \draw[arr] (tracks) -- (veto);
      \end{tikzpicture}
    \end{column}
    \begin{column}{0.43\textwidth}\footnotesize
      Built in \texttt{fast-veto-sim} (Python + numpy at run time), used as it is.\\[1.5mm]
      \begin{tabular}{@{}lr@{}}
        \toprule
        step & time\\
        \midrule
        muon interactions, one design (4\,M) & 60--80\,s\\
        cache: helium, PVC, air (11\,M, delta) & 89\,s\\
        neutrinos (4.2\,M, GENIE library) & $\sim$90\,s\\
        \key{one design, all backgrounds} & \key{3--6\,s}\\
        \bottomrule
      \end{tabular}\\[1.5mm]
      \takeaway[\footnotesize]{Helium, PVC and air do not depend on the design: drawn once. Each design only needs its veto.}
    \end{column}
  \end{columns}
\end{frame}

\begin{frame}{Generators: Geant4, Pythia\,6 and GENIE}
  \begin{columns}[T]
    \begin{column}{0.5\textwidth}\footnotesize
      \key{Libraries}, built once per material and energy (fast-veto-sim):
      \begin{itemize}\scriptsize
        \item \key{muons}: Geant4 muonNuclear for the soft part (resonances, nuclear break-up), \key{Pythia\,6} for the hard part ($W > 2$\,GeV); our default \emph{crossover}: Geant4 below an energy transfer $\nu_x(E)$, Pythia above. 50 energies 0.5--500\,GeV, 20\,k events per energy and material.
        \item \key{neutrinos}: \key{GENIE 3.06.02} (G18\_02a, Pythia\,6 hadronisation, charm and $\tau$ decays): $\nu_\mu$, $\bar\nu_\mu$, $\nu_e$, $\bar\nu_e$ on every nucleus of our materials, 1--500\,GeV.
        \item short-lived products (K$^0_S$, hyperons, $\pi^0$, $\eta$) decayed in the library, decay point from $c\tau$.
      \end{itemize}
    \end{column}
    \begin{column}{0.47\textwidth}\footnotesize
      \key{Applying} a library to our muons (or the neutrino flux):
      \begin{enumerate}\scriptsize
        \item every crossing of helium, PVC, air and each plane: $P = n_\text{area}\,\sigma(E)$;
        \item per target, \key{force} $M$ interactions drawn $\propto wP$, each of weight $S/M$, $S = \sum wP$: the weights add up to the true rate;
        \item vertex uniform along the crossing; a library event at a neighbouring energy, rotated onto the direction, random azimuth;
        \item the scattered muon continues from its \key{parent's own path}: the parent's real hits upstream are digitised.
      \end{enumerate}
      \takeaway[\footnotesize]{Helium, PVC and air do not depend on the design: drawn once and cached. A design's own planes are drawn per design.}
    \end{column}
  \end{columns}
\end{frame}

\begin{frame}{Forced interactions}
  \begin{columns}[c]
    \begin{column}{0.52\textwidth}\small
      Real interactions are $10^{-6}$--$10^{-8}$ per crossing: far too rare to simulate muon by muon. So:
      \begin{enumerate}\footnotesize
        \item all crossings, $P = n_\mathrm{area}\sigma(E)$; per target group $g$ the total $S_g = \sum wP$ (interactions per spill);
        \item draw $M_g$ crossings with probability $\propto wP$ and \key{make them interact}, each with weight $S_g/M_g$: the weights add up to the true rate;
        \item vertex uniform along the crossing;
        \item final state: a library event at a neighbouring energy, rotated onto the muon's direction, random azimuth.
      \end{enumerate}
      \aside{Every number later is $\sum$ weight $\times\,1.2\cdot10^{7}$ spills (15 years); errors: bootstrap over input muons.}
    \end{column}
    \begin{column}{0.45\textwidth}\footnotesize
      Muon interactions per 15 years:\\[1mm]
      \begin{tabular}{@{}lr@{}}
        \toprule
        helium & $1.5\cdot10^{6}$\\
        PVC walls & $2.2\cdot10^{7}$\\
        UBT (straw) & $3.2\cdot10^{4}$\\
        one tile fin at 30.6\,m & $\sim10^{11}$\\
        cavern air, 30.6--82\,m & $1.4\cdot10^{12}$\\
        \bottomrule
      \end{tabular}\\[1.5mm]
      Helium agrees with the full-Geant4 rate study (1.46\,M).\\[1.5mm]
      \takeaway[\footnotesize]{Most happen outside the vessel (air, fins): their products rarely make a vertex inside. The helium and the walls matter.}
    \end{column}
  \end{columns}
\end{frame}

\begin{frame}{Cross sections}
  \livewidget{sim_xs_mu}{\linewidth}{0.6\textheight}\\[1mm]
  \begin{columns}[c]
    \begin{column}{0.4\textwidth}\livewidget{sim_xs_split}{\linewidth}{0.62cm}\end{column}
    \begin{column}{0.56\textwidth}\livewidget{sim_xs_mat}{\linewidth}{0.62cm}\end{column}
  \end{columns}
  \legendline{crossover: Geant4 below $\nu_x(E)$, Pythia above · w2: Geant4 below $W = 2$\,GeV · geant4: Geant4 only. Pythia gives 1.4--1.7$\times$ more extra tracks: the generator systematic.}
\end{frame}

\begin{frame}{Grazing walls: a factor 4.4}
  \begin{columns}[c]
    \begin{column}{0.66\textwidth}
      \livewidget{sim_wall}{\linewidth}{0.62\textheight}\\[1mm]
      \sld{sim_w_ang}
    \end{column}
    \begin{column}{0.32\textwidth}\footnotesize
      The side walls lean outward by 30\,mrad. Muons cross them at a few degrees or less, where the path in 0.3\,mm of PVC
      is centimetres to metres (here: 40\,mm on average).\\[1.5mm]
      Our old scattering capped it at $50\times$ the thickness: \key{4.9\,M} interactions in the walls per 15 years. Exact chords:
      \key{21.6\,M}, as the full-Geant4 studies (18--19\,M). Now in our scattering too.\\[1.5mm]
      \aside{Right: angle to the wall of muons crossing a side wall; red line: the old cap.}
    \end{column}
  \end{columns}
\end{frame}

\begin{frame}{Where the interactions are}
  \begin{columns}[c]
    \begin{column}{0.7\textwidth}
      \livewidget{sim_vtx3d}{\linewidth}{0.58\textheight}\\[1mm]
      \livewidget{sim_v_tgt}{\linewidth}{0.62cm}\\[0.5mm]
      \livewidget{sim_v_kind}{0.5\linewidth}{0.62cm}
    \end{column}
    \begin{column}{0.28\textwidth}\footnotesize
      Light blue: helium vertices, red: wall vertices (samples); black: the chosen target.\\[2mm]
      Muon vertices in the helium follow the muons that enter the vessel; neutrino vertices fill it around the axis.\\[2mm]
      \aside{D6\_straw; the helium, PVC and air rates are the same for every design.}
    \end{column}
  \end{columns}
\end{frame}

\begin{frame}{Neutrinos}
  \begin{columns}[c]
    \begin{column}{0.5\textwidth}\small
      fast-veto-sim's own \key{GENIE 3.06.02} (G18\_02a, Pythia\,6 hadronisation, charm and $\tau$ decays): a library of
      $\nu_\mu$, $\bar\nu_\mu$, $\nu_e$, $\bar\nu_e$ on every nucleus of our materials, 1--500\,GeV, applied to the
      neutrino flux through the same forced-interaction machinery.\\[1.5mm]
      \key{No incoming charged particle}: nothing for a veto plane to see. Only the selection can remove them.\\[1.5mm]
      \aside{One neutrino interaction at a time: Part 2, \emph{One applied interaction at a time}, from: neutrinos.}
    \end{column}
    \begin{column}{0.47\textwidth}\footnotesize
      Per 15 years (D6\_straw; helium, PVC and air are design independent):\\[1mm]
      \begin{tabular}{@{}lrr@{}}
        \toprule
        target & interactions & $\geq2$ tracks in SST1\\
        \midrule
        helium & 5\,290 & 1\,360\\
        PVC walls & 3\,690 & 374\\
        air & 5\,690 & 881\\
        \bottomrule
      \end{tabular}\\[1.5mm]
      After our selection (Part 5): \key{0.8} per 15 years, $\approx0.5$ after the veto's random hits: a floor
      for every design.
    \end{column}
  \end{columns}
\end{frame}

\begin{frame}{Multiple scattering: the tails}
  \begin{columns}[c]
    \begin{column}{0.64\textwidth}
      \livewidget{sim_msc_hist}{\linewidth}{0.58\textheight}\\[1mm]
      \sld{sim_ms_p}\\[0.3mm]
      \livewidget{sim_ms_q}{0.5\linewidth}{0.62cm}
    \end{column}
    \begin{column}{0.34\textwidth}\footnotesize
      Our scattering (now with fast-veto-sim's tails): Highland's Gaussian core over walls, planes, \key{helium and air}
      (1\,m slices, on now), exact wall chords.\\[1.5mm]
      \key{Tails}: rare large single scatters. Above $2\,\theta_0$, a Poisson number per scatterer from Rutherford,
      $1/\theta^4$ up to the nuclear size; one parameter tuned to Geant4.\\[1.5mm]
      RMS and 99\,\% quantile at 82\,m: within 5--8\,\% of Geant4 at 1.5--100\,GeV; Gaussian alone: 99\,\% 25\,\% too small.
    \end{column}
  \end{columns}
\end{frame}

\begin{frame}{Fast chain vs full Geant4}
  \begin{columns}[c]
    \begin{column}{0.66\textwidth}
      \livewidget{sim_closure}{\linewidth}{0.66\textheight}
    \end{column}
    \begin{column}{0.32\textwidth}\footnotesize
      2\,M interactions in helium and PVC in full Geant4 (everything transported to 82\,m) and the same draws through the
      fast chain. Charged tracks besides the muon, $p > 1$\,GeV, in SST1:\\[1mm]
      \begin{tabular}{@{}lrr@{}}
        fast vs full (helium) & $\geq1$ & $\geq2$\\\midrule
        Geant4 physics & $-5\,\%$ & $-9\,\%$\\
        crossover (ours) & $\times1.3$ & $\times1.6$\\
      \end{tabular}\\[1.5mm]
      Missing in the fast chain: photon conversions, secondary interactions. Our default (crossover) is the
      \key{conservative} side for two-track vertices.
    \end{column}
  \end{columns}
\end{frame}

% ======================================================================== PART 4
\blurb{A UBT on the vessel entrance and thin fins reaching out from the side walls, left and right. Nothing on top or below.}
\section{The design space}

\begin{frame}{UBT and fins}
  \livewidget{design3d}{0.63\linewidth}{0.46\textheight}\hfill\livewidget{design_eff}{0.35\linewidth}{0.46\textheight}\\[1mm]
  \sld[0.32\linewidth]{d_ubt}\hfill\sld[0.32\linewidth]{d_nfin}\hfill\sld[0.32\linewidth]{d_gap}\\[0.4mm]
  \sld[0.32\linewidth]{d_len}\hfill\sld[0.32\linewidth]{d_top}\hfill\hspace*{0.32\linewidth}\\[0.4mm]
  \legendline{Sliders: UBT top · fins per side · gap · fin length · fin top [m]. Fin colour: $\log_{10}$ muons/m$^2$/spill, 2 (purple) to 9 (yellow). Black: drawn muons crossing nothing. Bars: geometric coverage.}
\end{frame}

\begin{frame}{Where the muons are}
  \begin{columns}[c]
    \begin{column}{0.6\textwidth}
      \livewidget{rate_map}{\linewidth}{0.66\textheight}\\[1mm]
      \sld[0.48\linewidth]{map_z}\hfill\sld[0.48\linewidth]{map_side}
    \end{column}
    \begin{column}{0.37\textwidth}\footnotesize
      Muon density on a fin plane at $z$, every muon of the dataset; $u$ is the distance from the vessel wall. Side: 0 right ($-x$), 1 left ($+x$). Dotted: cavern wall and floor limit.\\[2mm]
      A \key{lobe} at $y\approx0$, 1.5--3\,m from the wall, reaches $10^{9}$ muons/m$^2$ per spill: the muons the shield sweeps past the vessel. Close to the wall and away from $y=0$ it is 3--5 orders of magnitude quieter.\\[2mm]
      \takeaway[\footnotesize]{A fin pays for every muon crossing it. The muons it is there for cross it near the wall.}
    \end{column}
  \end{columns}
\end{frame}

\begin{frame}{Seven designs to learn from}
  \begin{columns}[c]
    \begin{column}{0.55\textwidth}\footnotesize
      \begin{tabular}{@{}lp{0.62\linewidth}r@{}}
        \toprule
        & & area\\
        \midrule
        D0 & UBT on the entrance only & 3\,m$^2$\\
        D1 & taller UBT, one 4\,m fin per side at 30.6\,m & 56\,m$^2$\\
        D2 & 4 fins per side every 15\,m, 3\,m long, $y\pm3$\,m & 145\,m$^2$\\
        D3 & 7 fins on the right up to the wall, 3 short on the left & 209\,m$^2$\\
        D4 & 8 fins per side, 5\,m, $y$ to 6\,m: the stress test & 624\,m$^2$\\
        D5 & D2 with 50\,cm or 5\,cm cells & 145\,m$^2$\\
        D6 & D2-like with real technologies: straw, fibre, pixel, tile, tile2 & --\\
        E1--E3 & evolved: a minimal UBT and two short tile2 fin pairs (30.6, 36\,m) & 4\,m$^2$\\
        \bottomrule
      \end{tabular}
    \end{column}
    \begin{column}{0.42\textwidth}\footnotesize
      A design is a small JSON file: UBT edges and cell size; per fin its side, $z$, length, $y$ range, cell size; detector efficiency and time resolution.\\[2mm]
      The loader applies the physical limits (floor, cavern wall, minimum UBT) and says what it changed.\\[2mm]
      All: 10\,cm cells, $\varepsilon = 0.99$, $\sigma_t = 0.5$\,ns, no tilt yet.\\[2mm]
      \takeaway[\footnotesize]{In the RL environment fins are placed one after another along $z$ with a minimum gap.}
    \end{column}
  \end{columns}
\end{frame}

% ======================================================================== PART 5
\blurb{Events are short time windows around a spectrometer candidate. The selection keeps two-track vertices pointing at the target; every track is followed back to the veto planes, in a window that grows from the vertex.}
\section{Reconstruction}

\begin{frame}{The chain}
  \centering
  \begin{tikzpicture}[node distance=4mm and 5mm]
    \node[hbox, text width=2.7cm] (paths) {\semibold true paths\\\faint{muons A $\to$ K $\to$ B; interaction products}};
    \node[box, right=of paths, text width=2.7cm] (dig) {\semibold digitization\\\faint{crossing $\to$ cell, $\varepsilon$, time $\pm\sigma_t$}};
    \node[box, below=of dig, text width=2.7cm] (acc) {\semibold accidental hits\\\faint{Poisson(rate $\times$ window) per plane, from rate maps}};
    \node[abox, right=of dig, text width=2.0cm, yshift=-6mm] (hits) {\semibold hits\\\faint{plane, cell, time}};
    \node[hbox, below=of paths, text width=2.7cm, yshift=-12mm] (spec) {\semibold spectrometer tracks\\\faint{position, slope, time, $p$ smeared}};
    \node[box, right=of spec, text width=2.7cm, yshift=-10mm] (vtx) {\semibold selection\\\faint{two tracks, DOCA $<2$\,cm, measured IP $<0.1$\,m, no third track}};
    \node[abox, right=of hits, text width=2.6cm, yshift=-12mm, line width=1.2pt] (algo) {\semibold the veto\\\faint{every track back, window + vertex cone, time}};
    \node[box, below=of algo, text width=2.6cm] (score) {\semibold veto / no veto\\\faint{pairs, $\mu$ and $\nu$ interactions\\signal: FP · TN}};
    \draw[arr] (paths) -- (dig); \draw[arr] (dig) -- (hits); \draw[arr] (acc) -- (hits);
    \draw[arr] (paths) -- (spec); \draw[arr] (spec) -- (vtx); \draw[aarr] (hits) -- (algo);
    \draw[aarr] (vtx.east) -- node[below, font=\scriptsize, color=accent] {tracks + vertex} ++(1.2,0) |- ([yshift=-2mm]algo.west);
    \draw[aarr] (algo) -- (score);
  \end{tikzpicture}\\[3mm]
  \footnotesize Four samples, one chain: muon pairs, muon and neutrino interactions, and the signal (no own veto hits: only random hits veto it).
\end{frame}

\begin{frame}{Events are time windows, not spills}
  \begin{columns}[c]
    \begin{column}{0.62\textwidth}
      \livewidget{acc_cloud}{\linewidth}{0.62\textheight}\\[1mm]
      \sld[0.6\linewidth]{cell_cm}\\[0.3mm]
      \legendline{cell size [cm]; 1\,s spill}
    \end{column}
    \begin{column}{0.36\textwidth}\footnotesize
      A fin at 30.6\,m sees $\sim10^{9}$ muons per second: simulating whole spills hit by hit is hopeless.\\[1.5mm]
      Instead, around each candidate: the own muons' hits, plus per plane
      \[ n_\text{acc} \sim \text{Poisson}\Big(R_\text{plane}\cdot\frac{2W}{T_\text{spill}}\Big) \]
      random hits, placed according to the full-statistics rate map and digitized like real ones.\\[1.5mm]
      Blue box: a search window near the wall. What matters is the chance a random hit lands in it: rate density $\times$ box $\times$ window.\\[1.5mm]
      \legendline{Orange squares: random hits of one $\pm1.5$\,ns window on the right fin of D2, at the chosen cell size.}
    \end{column}
  \end{columns}
\end{frame}

\begin{frame}{The detector and spectrometer models}
  \begin{columns}[T]
    \begin{column}{0.48\textwidth}\footnotesize
      \key{Veto planes} (UBT and every fin, the same model)
      \begin{itemize}\scriptsize
        \item one layer of cells at fixed $z$ (10\,cm default)
        \item a crossing makes one hit with probability $\varepsilon$ (0.99)
        \item hit = cell centre + time, $\sigma_t = 0.5$\,ns, muons at $c$
        \item no direction, no energy, no noise, no dead time (yet)
        \item a muon's own crossing comes from its true path, bounce included
      \end{itemize}
      \vspace{1mm}
      \key{Spectrometer} (track at $z = 82$\,m)
      \begin{itemize}\footnotesize
        \item position $\pm1$\,mm, time $\pm0.1$\,ns
        \item slope $\pm0.1$\,mrad $\oplus$ 2\,mrad$\cdot$(1\,GeV$/p$)
        \item momentum as true, used only to size the window
      \end{itemize}
    \end{column}
    \begin{column}{0.48\textwidth}\footnotesize
      \key{Direction comes only from the spectrometer.} The UBT gives no usable direction; neither do the fins.\\[1.5mm]
      \key{Samples}
      \begin{itemize}\scriptsize
        \item combinatorial: opposite-charge muon pairs of one seed in the selection, weight $w_1 w_2 N_\text{runs}\cdot 2\Delta t/T_\text{spill}$, $\Delta t = 200$\,ps: 197\,k (alpha); 25\,M (delta), subsampled $\propto$ weight to 400\,k (total within 0.05\,\%)
        \item signal: long-lived neutrals from the target ($\theta\approx p_T/p$), decaying in the vessel into two tracks
        \item single SST1 muons (200\,k) for per-muon efficiencies
        \item forced $\mu$ (11\,M) and $\nu$ (7--11\,M) interactions
      \end{itemize}
    \end{column}
  \end{columns}
\end{frame}

\begin{frame}{The signal: from the target, not from Mars}
  \begin{columns}[c]
    \begin{column}{0.56\textwidth}
      \livewidget{sig_cross}{\linewidth}{0.7\textheight}
    \end{column}
    \begin{column}{0.41\textwidth}\footnotesize
      A long-lived neutral from the \key{target}: $p$ = 10--100\,GeV (log-uniform), $p_T\sim$ exp(1\,GeV), so $\theta\approx p_T/p$.\\[1mm]
      Decays uniformly in $z$ in the vessel, $m = 1$\,GeV $\to$ two muon-mass tracks.\\[1mm]
      The pairs' selection: SST1, $p>1$\,GeV, timing, IP $<2.5$\,m. \key{37\,\%} pass; median daughter angle 22\,mrad.\\[1mm]
      Orange: its tracks followed back to $z = 30.6$\,m; grey: muon density.\\[1mm]
      \takeaway[\footnotesize]{Two thirds cross the quiet UBT; the rest hit the fins next to the wall, where random hits veto them.}
    \end{column}
  \end{columns}
\end{frame}

\begin{frame}{The selection}
  \begin{columns}[c]
    \begin{column}{0.5\textwidth}\small
      For \key{every} sample (signal, pairs, interactions):\\[1mm]
      \begin{enumerate}\footnotesize
        \item two opposite-charge tracks, $p > 1$\,GeV, in the SST1 and timing acceptance;
        \item DOCA $< 2$\,cm, vertex in the vessel;
        \item \key{pointing}: the measured tracks (smeared at 82\,m) give an IP at the target $< 0.1$\,m;
        \item \key{two tracks}: no other track in SST1 in the event.
      \end{enumerate}
      A signal decay from the target points back (measured IP: median 2\,cm, 99.9\,\% below 8.3\,cm) and has two
      tracks; interactions have more tracks, and their muons were swept by the shield.
    \end{column}
    \begin{column}{0.47\textwidth}\footnotesize
      What survives the new cuts (3 + 4):\\[1mm]
      \begin{tabular}{@{}lrr@{}}
        \toprule
        sample & before & after\\
        \midrule
        signal & 100\,\% & \key{99.98\,\%}\\
        combinatorial pairs & 36 & 2.5\\
        muon interactions & $1.25\cdot10^{5}$ & $2.0\cdot10^{3}$\\
        neutrino interactions & 481 & \key{0.8}\\
        \bottomrule
      \end{tabular}\\[1mm]
      \aside{Per 15 years, before the veto. Two-body signal model: decays with an invisible particle point back less well.}
    \end{column}
  \end{columns}
\end{frame}

\begin{frame}{Algorithm A: follow every track back}
  \begin{columns}[c]
    \begin{column}{0.48\textwidth}\footnotesize
      For each spectrometer track:
      \begin{enumerate}\footnotesize
        \item extrapolate backward in a straight line from $z = 82$\,m (no field in the vessel)
        \item on every veto plane \key{upstream of the vertex}: predicted cell $(u, y)$ and time $t - \ell/c$
        \item look for a hit within
          \[ |\Delta u|, |\Delta y| \le r_0 + \Big(a + \frac{b}{p}\Big)d + \tfrac{\text{cell}}{2},\quad |\Delta t| \le n\,\sigma \]
        \item a hit found: the track is \key{explained}
      \end{enumerate}
      A candidate is vetoed if either track is explained.\\[2mm]
      \takeaway[\footnotesize]{No combinatorics: one binary search per (event, plane) and a scan of that plane's few hits.}
    \end{column}
    \begin{column}{0.5\textwidth}\centering
      \begin{tikzpicture}[x=0.95cm,y=0.95cm]
        \fill[wash] (0,-0.8) -- (5.6,-1.7) -- (5.6,1.7) -- (0,0.8) -- cycle;
        \draw[accent] (0,-0.8) -- (5.6,-1.7) (0,0.8) -- (5.6,1.7);
        \draw[grn, line width=2pt] (1.2,-0.99) -- (1.2,-2.9); \node[font=\scriptsize, color=grn, left] at (1.2,-2.8) {fin};
        \draw[ora, line width=2pt] (0,-0.8) -- (0,0.8); \node[font=\scriptsize, color=ora, left] at (0,0.6) {UBT};
        \draw[red, line width=1.2pt] (-0.6,-2.6) -- (2.6,-1.25) -- (6.2,0.3);
        \draw[red, dashed] (2.6,-1.25) -- (-0.6,-2.63);
        \fill[accent, opacity=0.25] (1.08,-2.35) rectangle (1.32,-1.7);
        \fill[white, draw=red, line width=1pt] (1.2,-2.02) circle (0.07);
        \node[font=\scriptsize, color=accent, right] at (1.35,-2.45) {window};
        \fill[red] (6.2,0.3) circle (0.06); \node[font=\scriptsize, right] at (6.2,0.3) {track};
        \draw[muted] (6.2,-1.9) -- (6.2,1.9); \node[font=\scriptsize, color=muted, above] at (6.2,1.9) {SST1};
        \fill[red] (4.3,-0.5) circle (0.06); \node[font=\scriptsize, color=red, below] at (4.3,-0.55) {vertex};
        \draw[muted, dotted] (4.3,-2.9) -- (4.3,2.0); \node[font=\scriptsize, color=muted] at (4.3,2.2) {only upstream};
      \end{tikzpicture}
    \end{column}
  \end{columns}
\end{frame}

\begin{frame}{The veto algorithm, step by step}
  \begin{columns}[c]
    \begin{column}{0.48\textwidth}\centering
      \begin{tikzpicture}[node distance=2.4mm, every node/.append style={font=\scriptsize}]
        \node[hbox, text width=5cm] (t) {\semibold every track of the event in SST1\\\faint{measured: position, slopes, time}};
        \node[box, below=of t, text width=5cm] (p) {\semibold straight back to every plane upstream of the vertex\\\faint{UBT and fins, $z_\mathrm{plane} < z_\mathrm{vertex}$}};
        \node[abox, below=of p, text width=5cm] (w) {\semibold search window on the plane\\\faint{$r = \max\big(r_A,\ r_0 + \theta_\mathrm{cone}\,(z_\mathrm{v} - z_\mathrm{plane})\big)$}};
        \node[box, below=of w, text width=5cm] (m) {\semibold a hit inside $r$ and in time?\\\faint{$\pm3\sigma_t$; tiles $\pm1.5$\,ns; straws 50\,ns drift}};
        \node[hbox, below=of m, text width=5cm] (v) {\semibold yes on any track: event vetoed};
        \draw[arr] (t) -- (p); \draw[aarr] (p) -- (w); \draw[arr] (w) -- (m); \draw[arr] (m) -- (v);
      \end{tikzpicture}
    \end{column}
    \begin{column}{0.49\textwidth}\scriptsize
      \key{Algorithm A} (as before): $r_A = r_0 + (r_\mathrm{slope} + r_\mathrm{ms}/p)\,d + \tfrac{\text{cell}}{2}$, $d$ the
      distance from SST1: $r_0 = 5$\,cm, $r_\mathrm{slope} = 0.5$\,mrad, $r_\mathrm{ms} = 6$\,mrad\,GeV. Hits: the event's own
      (efficiency, cells, timing) and random ones from the rate maps.\\[1.5mm]
      \key{New 1, every track}: not only the two vertex tracks. The scattered muon of an interaction, followed back
      past the vertex, points at its parent's hit.\\[1.5mm]
      \key{New 2, the vertex cone}: the parent was bent at the vertex by a few mrad, so its hit is off the backward line
      by $\approx$ angle $\times$ distance \emph{from the vertex}: the window grows from the vertex.
      A design parameter: wider catches more parents and more random hits (signal loss).\\[1mm]
      \faint{Single muons: no vertex, unchanged.}
    \end{column}
  \end{columns}
\end{frame}

\begin{frame}{The spatial window}
  \begin{columns}[c]
    \begin{column}{0.6\textwidth}
      \livewidget{window_plot}{\linewidth}{0.64\textheight}\\[1mm]
      \sld[0.6\linewidth]{r_ms}
    \end{column}
    \begin{column}{0.37\textwidth}\footnotesize
      $r = r_0 + (a + b/p)\,d$ with $r_0 = 5$\,cm, $a = 0.5$\,mrad and $b$ the slider, in mrad$\cdot$GeV: the spectrometer's slope error grows with the lever arm $d$, multiple scattering with $1/p$.\\[2mm]
      Title: how many of the true crossings of straight SST1 muons (D2, 60\,k muons) fall inside the window, computed now.\\[2mm]
      \takeaway[\footnotesize]{For straight muons the prediction is within millimetres: the window can be tight. A bigger window only buys more random hits.}
    \end{column}
  \end{columns}
\end{frame}

\begin{frame}{The time window}
  \livewidget{time_plot}{0.55\linewidth}{0.5\textheight}\hfill\livewidget{time_trade}{0.43\linewidth}{0.5\textheight}\\[1mm]
  \sld[0.48\linewidth]{sig_t}\hfill\sld[0.48\linewidth]{n_sig}\\[1mm]
  \footnotesize The predicted time includes the time of flight from the plane to $z = 82$\,m. The window is $\pm n\sigma$ with $\sigma = \sigma_t^\text{veto}\oplus\sigma_t^\text{track}$: wider loses fewer own hits, but lets in proportionally more random hits.\\
  \legendline{Right: signal loss scaled from D2's 1.3\,\% at $3\sigma$ and 0.5\,ns, since random hits grow with the window.}
\end{frame}

\begin{frame}{Why the parent is missed, and the cone}
  \livewidget{iv_event3d}{\linewidth}{0.6\textheight}\\[1mm]
  \begin{columns}[c]
    \begin{column}{0.44\textwidth}\sld{iv_ev}\end{column}
    \begin{column}{0.54\textwidth}\livewidget{iv_cone}{\linewidth}{0.62cm}\end{column}
  \end{columns}
  \legendline{Selected muon interactions (helium / walls), E3\_robust, the heaviest 200. Black: the parent up to the vertex; red: the scattered muon; blue: products. Switch the cone and watch the windows grow.}
\end{frame}

\begin{frame}{Algorithm A on real events}
  \livewidget{ev3d}{\linewidth}{0.52\textheight}\\[1mm]
  \sld[0.32\linewidth]{ev_design}\hfill\sld[0.32\linewidth]{ev_kind}\hfill\sld[0.32\linewidth]{ev_index}\\[0.3mm]
  \legendline{design: 0 D0 · 1 D1 · 2 D2 · 3 D3 · 4 D4 · 5 D5 coarse · 6 D5 fine · 7 fibre · 8 pixel · 9 straw · 10 tile · 11 tile2 (fins) · 12 E1 · 13 E2 (LLM).\quad kind: 0 pair vetoed · 1 pair missed · 2 signal lost to a random hit · 3 bounced muon missed.}\\
  \legendline{Solid: true path; dotted: backward prediction; white: own hits; diamond: the hit that explained the track; grey: random hits.}
\end{frame}

\begin{frame}{The same event from above}
  \livewidget{ev_side}{\linewidth}{0.55\textheight}\\[1mm]
  \sld[0.32\linewidth]{ev_design}\hfill\sld[0.32\linewidth]{ev_kind}\hfill\sld[0.32\linewidth]{ev_index}\\[0.3mm]
  \legendline{$x$--$z$ view. Thick lines: UBT (orange) and fins (green). Blue bars: the search window on each plane the track is followed to. Same sliders as the previous slide.}
\end{frame}

\begin{frame}{What it costs}
  \begin{columns}[c]
    \begin{column}{0.58\textwidth}\footnotesize
      \begin{tabular}{@{}lrrrr@{}}
        \toprule
        design & planes & hits / event & comparisons & µs / event\\
        \midrule
        D0 & 1 & 0.1 & 40\,k & 0.5\\
        D1 & 3 & 16 & 3.2\,M & 4.8\\
        D2 & 9 & 39 & 4.0\,M & 12.5\\
        D3 & 11 & 49 & 4.9\,M & 16.7\\
        D4 & 17 & 90 & 7.2\,M & 33.3\\
        \bottomrule
      \end{tabular}\\[1mm]
      \faint{196\,736 background events; numpy, one core of a laptop.}
    \end{column}
    \begin{column}{0.4\textwidth}\small
      A full evaluation of a design (196\,k pairs, 100\,k signal, 200\,k single muons) takes \key{$\approx$6\,s} on one core, linear in hits per event.\\[2mm]
      Nothing to cap: algorithm A never combines hits.\\[2mm]
      The interaction backgrounds add 3--6\,s: only their veto runs per design.\\[2mm]\takeaway[\footnotesize]{Tens of seconds per design on the cluster, 8 in parallel: fast enough for an evolution loop.}
    \end{column}
  \end{columns}
\end{frame}

% ======================================================================== PART 6
\blurb{One design through everything: the enriched dataset, every background source before and after the veto, the vertex cone, the signal, the cost.}
\section{Results}

% End to end on the enriched dataset (bin/end_to_end.py, results/end_to_end/e2e.json); widgets in talk_e2e.py.

\begin{frame}{The designs, and two iterations}
  \begin{columns}[c]
    \begin{column}{0.55\textwidth}
      \livewidget{e2e_design3d}{\linewidth}{0.58\textheight}\\[1mm]
      \livewidget{e2e_design}{\linewidth}{0.62cm}
    \end{column}
    \begin{column}{0.43\textwidth}\footnotesize
      \key{E3\_robust}: from the first evolution runs (E2 + 10\,cm margins; see Evolution): the minimal straw UBT and two short tile2 fin pairs at 30.6 and 36\,m, 26\,k channels, 0.40\,MCHF. \key{D6\_tile2}: the hand-made reference, tile2 fins every 15\,m.\\[1.5mm]
      \key{Iteration 1}: E3's surviving muon interactions all have a parent that crossed \emph{one} plane, the UBT or a 30.6\,m fin, and left no hit: plane efficiency 0.999. Hence a \key{second station} where the parents enter:
      \begin{itemize}\scriptsize
        \item \key{E4\_ubt2}: two UBT stations ($+97$\,kCHF)
        \item \key{E4\_entry2}: two stations on the UBT and both 30.6\,m fins ($+206$\,kCHF)
      \end{itemize}
      \key{Iteration 2}: the cone of E4\_entry2, 0--20\,mrad.\\[1mm]
      \aside{A second station: efficiency $1 - 10^{-6}$ and its cost added; same geometry, so the same interactions.}
    \end{column}
  \end{columns}
\end{frame}

\begin{frame}{Every background, before and after the veto}
  \begin{columns}[c]
    \begin{column}{0.64\textwidth}
      \livewidget{e2e_sources}{\linewidth}{0.62\textheight}\\[1mm]
      \livewidget{e2e_cone}{\linewidth}{0.62cm}
    \end{column}
    \begin{column}{0.34\textwidth}
      \livewidget{e2e_design}{\linewidth}{0.62cm}\\[2mm]\footnotesize
      Light: selected, before the veto; dark: after it. Per 15 years, enriched dataset, the full chain: scattering with tails, helium and air; measured IP $<0.1$\,m; at most two tracks; the cone of the design.\\[1.5mm]
      Muon interactions in the PVC walls (130--170) and in the veto planes (75--80) are all vetoed: their parents crossed the planes right before. In the \key{helium}, 4\,300 are selected.\\[1.5mm]
      \takeaway[\footnotesize]{Neutrinos: $\approx$0.9 selected, a third vetoed by random hits; $\approx$0.6 is a floor for every design.}
    \end{column}
  \end{columns}
\end{frame}

\begin{frame}{Why the last muon interactions survive}
  \begin{columns}[c]
    \begin{column}{0.5\textwidth}\footnotesize
      The survivors of E3\_robust at 10\,mrad, one by one (the parent at $z = 30.6$\,m):\\[1mm]
      \begin{itemize}\scriptsize
        \item all from the \key{high-energy tail}: 130--300\,GeV, the enriched input's own muons
        \item parents enter at $|x| = 0.43$--$0.54$\,m: through the UBT window or the $-x$ wall right behind the 30.6\,m fin
        \item in truth every parent crossed a plane; the event simply has \key{no hit} from it: one plane, efficiency 0.999
        \item the scattered muon is reconstructed in every case; its kink is 1--2\,mrad
      \end{itemize}
      \vspace{1mm}
      $4\,300 \times 10^{-3} \approx 4$ per 15 years: exactly what is left once the cone recovers the kinks.
    \end{column}
    \begin{column}{0.46\textwidth}\footnotesize
      Two separate effects, two separate cures:\\[1mm]
      {\scriptsize\begin{tabular}{@{}lll@{}}
        \toprule
        missed because & how often & cure\\
        \midrule
        bend at the vertex & most (no cone) & vertex cone\\
        no hit on its plane & $10^{-3}$ & second station\\
        \bottomrule
      \end{tabular}}\\[2mm]
      \takeaway[\footnotesize]{Redundancy is not about geometry here: the parents of interactions come in through the entrance, where one plane is the only chance.}
    \end{column}
  \end{columns}
\end{frame}

\begin{frame}{Signal against background}
  \begin{columns}[c]
    \begin{column}{0.6\textwidth}
      \livewidget{e2e_trade}{\linewidth}{0.66\textheight}
    \end{column}
    \begin{column}{0.38\textwidth}\footnotesize
      Each line one design, each point a cone (labels in mrad); all background sources together, per 15 years, against the signal lost to random hits.\\[1.5mm]
      The cone costs signal through the random hits it lets in: cheap on short fins near the wall (E3, E4: 1\,\% at 10\,mrad), expensive on D6\_tile2's busy fins (6\,\% at 10, 18\,\% at 20\,mrad).\\[1.5mm]
      \takeaway[\footnotesize]{With a second station at the entrance the muon interactions go to zero at 17.5\,mrad: what is left is the neutrino floor.}
    \end{column}
  \end{columns}
\end{frame}

\begin{frame}{End-to-end figure of merit}
  \begin{columns}[c]
    \begin{column}{0.6\textwidth}\tiny
      \setlength{\tabcolsep}{2.4pt}
      \begin{tabular}{@{}lrrrrrrrr@{}}
        \toprule
        design & cone & signal lost & comb. & $\mu$ helium & $\nu$ helium & \key{total} & MCHF & reward\\
        \midrule
        E3\_robust & 0 & 0.06\,\% & $9\cdot10^{-5}$ & $76\pm6$ & 0.72 & $76\pm6$ & 0.40 & $-100.4$\\
         & 10 & 0.92\,\% & $4\cdot10^{-5}$ & $3.1\pm1.3$ & 0.64 & $3.8\pm1.3$ & & $-98.9$\\
         & 20 & 3.60\,\% & $3\cdot10^{-5}$ & $1.6\pm0.9$ & 0.54 & $2.1\pm0.9$ & & $-91.7$\\
        E4\_ubt2 & 10 & 0.90\,\% & $4\cdot10^{-5}$ & $1.6\pm0.9$ & 0.67 & $2.2\pm0.9$ & 0.50 & $-90.4$\\
         & 17.5 & 2.76\,\% & $4\cdot10^{-5}$ & $0.5\pm0.5$ & 0.62 & $1.1\pm0.5$ & & $-71.0$\\
        E4\_entry2 & 10 & 0.90\,\% & $2\cdot10^{-5}$ & $1.0\pm0.7$ & 0.67 & $1.7\pm0.7$ & 0.61 & $-83.1$\\
         & \key{17.5} & \key{2.85\,\%} & $1\cdot10^{-5}$ & \key{0} & 0.62 & \key{$0.62\pm0.06$} & & \key{$-49.2$}\\
         & 20 & 3.78\,\% & $1\cdot10^{-5}$ & 0 & 0.61 & $0.61\pm0.06$ & & $-49.6$\\
        D6\_tile2 & 10 & 5.89\,\% & $3\cdot10^{-5}$ & $1.0\pm0.7$ & 0.43 & $1.5\pm0.7$ & 4.4 & rates\\
         & 20 & 17.8\,\% & $1\cdot10^{-5}$ & $1.0\pm0.7$ & 0.37 & $1.4\pm0.7$ & & rates\\
        \bottomrule
      \end{tabular}\\[1.5mm]
      {\color{muted}\lightfont\tiny Per 15 years after selection and veto; $\mu$ interactions in air, PVC and the planes and $\nu$ in air and PVC: 0 in every row. 100\,k signal decays, 400\,k pairs, 11\,M + 7\,M forced interactions; delta dataset, 484 seeds.}
    \end{column}
    \begin{column}{0.38\textwidth}\footnotesize
      Confusion matrix of \key{E4\_entry2 at 17.5\,mrad}, per 15 years:\\[1mm]
      {\scriptsize\setlength{\tabcolsep}{3pt}\begin{tabular}{@{}lrrr@{}}
        \toprule
        & before & TP & FN\\
        \midrule
        combinatorial & 16.8 & 16.8 & $10^{-5}$\\
        $\mu$ helium & 4\,307 & 4\,307 & 0\\
        $\mu$ PVC walls & 128 & 128 & 0\\
        $\mu$ veto planes & 80 & 80 & 0\\
        $\nu$ helium & 0.92 & 0.30 & 0.62\\
        $\nu$ veto planes & 0.06 & 0.06 & 0\\
        \midrule
        signal & & FP 2.85\,\% & TN 97.2\,\%\\
        \bottomrule
      \end{tabular}}\\[1.5mm]
      \takeaway[\scriptsize]{0 muon interactions: fewer than one simulated event of weight 0.52, i.e.\ below $\approx$1.2 per 15 years at 90\,\% CL.}
    \end{column}
  \end{columns}
\end{frame}



\begin{frame}{What the first round taught}
  \begin{columns}[c]
    \begin{column}{0.55\textwidth}\small
      The first round (alpha dataset, combinatorial pairs only, no interactions, simpler scattering) is superseded by the numbers above. Its lessons hold:
      \begin{itemize}\footnotesize
        \item \key{time resolution} is the biggest lever on signal loss (0.5 $\to$ 0.1\,ns: $3\times$ less), then cell size
        \item a fin pays for every muon crossing it: the muons it is there for cross it \key{near the wall}, not in the hot lobe 1.5--3\,m out
        \item a single plane misses $\approx$1\,\% of the muons it should see (efficiency, time tails): \key{redundancy} has to be in the geometry, a 2-plane rule fails
        \item bounced muons are the blind spot of a backward search (algorithm B)
      \end{itemize}
    \end{column}
    \begin{column}{0.4\textwidth}\footnotesize
      \takeaway[\footnotesize]{The end-to-end study adds one more: the parents of interactions come in through the entrance, where a single plane with efficiency 0.999 is their only chance.}
    \end{column}
  \end{columns}
\end{frame}

% ======================================================================== PART 7
\blurb{An LLM proposes complete designs, the cluster scores them on every background, the best ones breed.}
\section{Optimisation}

\begin{frame}{The loop: an LLM proposes, the cluster scores}
  \begin{columns}[c]
    \begin{column}{0.44\textwidth}\centering
      \begin{tikzpicture}[node distance=2.4mm, every node/.append style={font=\scriptsize}]
        \node[abox, text width=4.3cm] (llm) {\semibold an LLM, 8 in parallel\\\faint{Claude Code or OpenAI Codex, headless; a complete design JSON + rationale}};
        \node[box, below=of llm, text width=4.3cm] (eval) {\semibold the cluster, 8 cores\\\faint{pairs, 40\,k signal decays, 200\,k muons, $\mu$ and $\nu$ interactions}};
        \node[box, below=of eval, text width=4.3cm] (obj) {\semibold reward\\\faint{\texttt{veto/lib/veto/objective.py} (G.~Frisella)}};
        \node[hbox, below=of obj, text width=4.3cm] (pop) {\semibold population of 12\\\faint{truncation selection}};
        \draw[aarr] (llm) -- (eval); \draw[arr] (eval) -- (obj); \draw[arr] (obj) -- (pop);
        \draw[arr] (pop.west) -- ++(-0.35,0) |- (llm.west);
      \end{tikzpicture}
    \end{column}
    \begin{column}{0.53\textwidth}\footnotesize
      AlphaEvolve-style, with \key{designs instead of code}. Each call sees:
      \begin{itemize}\scriptsize
        \item \key{once}: geometry and limits, how the veto works, technologies, muon flux by side, $z$ and distance from the wall, the objective, the JSON schema
        \item \key{per child}: a directive (refine / explore / combine), the parent and 3 inspirations with their reports, the population's rewards
        \item \key{per design}: reward $\pm$ error, violations, tiers, false vertices, signal efficiency, cost; per plane: peak rate vs limit, cost, signal lost, pairs vetoed
      \end{itemize}
      At the end the top 5 are rerun with two more seeds: the best \key{mean} wins.\\[1mm]
      Backends: \texttt{--backend claude} (Claude Code CLI) or \texttt{codex} (\texttt{codex exec}), both on the user's own login, no tools, no session saved; usage limits are waited out.\\[1mm]\takeaway[\footnotesize]{Old runs: 25 + 15 generations, $\approx$340 LLM calls, 256 designs: an evening on one node.}
    \end{column}
  \end{columns}
\end{frame}

\begin{frame}{The objective}
  \begin{columns}[c]
    \begin{column}{0.55\textwidth}\footnotesize
      Reward $= -\sum$ tiers, lexicographic through the weights:\\[1.5mm]
      \setlength{\tabcolsep}{3pt}\scriptsize
      \begin{tabular}{@{}llr@{}}
        \toprule
        tier & term & weight\\
        \midrule
        rate limits & $1 + \log_{10}(\text{peak}/\text{limit})$ per plane over it & 1000\\
        budget & $1 + $ overshoot / 8\,MCHF & 1000\\
        background & $1 - e^{-N}$, $N$ = all false vertices in 15\,y & 100\\
        signal & signal lost / 10\,\% & 10\\
        cost & $\log(\text{cost}/97\,\text{kCHF}) / \log 82$ & 1\\
        \bottomrule
      \end{tabular}\\[2mm]\footnotesize
      Exchange rates: \key{0.01 false vertices} $=$ \key{1\,\% signal} $=$ a factor 82 in cost.\\[1mm]
      Every term carries its statistical error: rewards closer than $\approx 2\sigma \approx 0.1$ are a tie.
    \end{column}
    \begin{column}{0.42\textwidth}\footnotesize
      $N$ = combinatorial + $\mu$ + $\nu$ interactions after selection and veto; the vertex cone is a design parameter (0--50\,mrad).\\[1mm]
      Combinatorial $N$ two ways:\\[1mm]
      \key{direct}: weighted surviving pairs, 0--20 of 196\,k.\\[1mm]
      \key{factorised}: $N_0 \sum_{kl} f_{kl}\,\varepsilon_k\varepsilon_l$, $\varepsilon_k$ how often a muon of type $k$ in a pair escapes the veto (measured on the pair muons, not the single-muon sample: bounced muons of pairs escape twice as often).\\[1.5mm]
      The optimiser uses the \key{larger} of the two.\\[1.5mm]
      \takeaway[\footnotesize]{Every hand-made design with fins was infeasible: the tile fins at 30.6\,m see 2--10\,MHz per channel (limit 1), the right straw fin $6\cdot10^5$\,Hz/cm$^2$ (limit $5\cdot10^5$), fibre costs 48\,MCHF.}
    \end{column}
  \end{columns}
\end{frame}

% ======================================================================== PART 8
\blurb{Two old runs on the combinatorial background only (outdated), and a new run on the full objective that found the next hole.}
\section{Evolution}

\begin{frame}{Three runs}
  \begin{columns}[c]
    \begin{column}{0.55\textwidth}\small
      \key{Run 3} (19 October) is the current one: the full objective on the enriched dataset, all three backgrounds,
      the vertex cone as a design parameter. \key{OpenAI Codex} (gpt-6.1-sol, high effort) on the cluster,
      10 generations $\times$ 8 children: 80 designs, 74 feasible, 15 minutes.\\[2mm]
      \key{Runs 1 and 2} (15 October, Claude Sonnet) optimised an older objective: alpha dataset, combinatorial
      background only, no interactions, no cone. \textcolor[HTML]{B0352B}{Outdated}; their best, E2, became E3 of the
      results.
    \end{column}
    \begin{column}{0.4\textwidth}\footnotesize
      Run 3 starts from the hand-made designs and E1--E3. Its seeds already show the new picture: the UBT alone has
      \key{1\,700} muon interactions per 15 years after the veto, E2 without a cone 91.\\[1.5mm]
      \takeaway[\footnotesize]{The runs share the slider on the next slides: 0, 1 old · \key{2 the new run} (default).}
    \end{column}
  \end{columns}
\end{frame}

\begin{frame}{Evolution}
  \begin{columns}[c]
    \begin{column}{0.62\textwidth}
      \livewidget{evo_curve}{\linewidth}{0.57\textheight}\\[1mm]
      \sld[0.6\linewidth]{evo_run}\\[0.3mm]
      \legendline{each dot one design (grey: infeasible), orange: the best so far. Hover for names. run: 0, 1 old (combinatorial only) · 2 run 3, full objective.}
    \end{column}
    \begin{column}{0.36\textwidth}\footnotesize
      \key{Run 3} starts at $-100$: every seed lets at least one false vertex through in 15 years.\\[1mm]
      Generation 1 adds a 22\,mrad cone and entrance fins: muon interactions 3, signal lost 2\,\%; still $-99.8$, since $100(1-e^{-N})$ hardly moves above $N \approx 3$.\\[1mm]
      Then the reward only improves by \key{widening the cone} (40, then 50\,mrad) and adding planes: $-70$, $-68$, $-24$, $-17$.\\[1.5mm]
      \takeaway[\footnotesize]{Old runs 1 and 2: from $-100$ to $-0.39$ by shrinking the fins (E1, E2); combinatorial background only.}
    \end{column}
  \end{columns}
\end{frame}

\begin{frame}{Watch the design evolve}
  \begin{columns}[c]
    \begin{column}{0.44\textwidth}
      \livewidget{evo3d}{\linewidth}{0.56\textheight}
    \end{column}
    \begin{column}{0.54\textwidth}
      \livewidget{evo_top}{\linewidth}{0.56\textheight}
    \end{column}
  \end{columns}
  \vspace{1mm}
  \sld[0.32\linewidth]{evo_run}\hfill\sld[0.62\linewidth]{evo_step}\\[0.3mm]
  \legendline{run: 0, 1 old · 2 run 3. Each step: a new best design. 3D: planes coloured by muon density. Top view: fins green / orange / red below 50\,\%, below 100\,\%, above their rate limit.}
\end{frame}

\begin{frame}{What the optimiser bought}
  \begin{columns}[c]
    \begin{column}{0.52\textwidth}
      \livewidget{evo_drift}{\linewidth}{0.7\textheight}
    \end{column}
    \begin{column}{0.46\textwidth}\tiny
      \setlength{\tabcolsep}{2pt}
      \begin{tabular}{@{}lrrrrr@{}}
        \toprule
        best at & reward & MCHF & sig.\ lost & $\mu$ & $\nu$\\
        \midrule
        seed (UBT) & $-100$ & 0.10 & 0 & 1\,724 & 0.58\\
        gen 1 & $-99.8$ & 0.70 & 2.1\,\% & 3.1 & 0.48\\
        gen 3 & $-70.5$ & 2.3 & 16\,\% & 0.5 & 0.39\\
        gen 5 & $-68.3$ & 7.2 & 36\,\% & 0.5 & 0.33\\
        gen 8 & $-23.8$ & 7.7 & 63\,\% & 0 & 0.14\\
        gen 10 & \key{$-17.0$} & 7.2 & \key{74\,\%} & 0 & \key{0.06}\\
        \midrule
        E4\_entry2 & $-49.2$ & 0.61 & 2.9\,\% & 0 & 0.62\\
        \bottomrule
      \end{tabular}\\[2mm]\footnotesize
      Every step trades signal for background: the muon interactions are gone by generation 8; after that the wider cone
      only removes \key{neutrinos}, by random hits, exactly as it removes the signal.
    \end{column}
  \end{columns}
\end{frame}

\begin{frame}{Where the evolved designs sit}
  \begin{columns}[c]
    \begin{column}{0.58\textwidth}
      \livewidget{evo_map}{\linewidth}{0.66\textheight}
    \end{column}
    \begin{column}{0.4\textwidth}\footnotesize
      Every feasible design of the three runs: cost against signal lost, coloured by all false vertices in 15 years.\\[1.5mm]
      Runs 1 and 2 (combinatorial only) sit at 0.3--0.5\,MCHF and 0.05--0.3\,\% signal lost. Run 3 climbs straight up the
      plot: \key{above the dotted line} (10\,\%) the signal term no longer grows, so more cone, more planes and more
      signal lost cost nothing extra.\\[1.5mm]
      \takeaway[\footnotesize]{The reasonable region, 1--3\,\% signal lost at under 1\,MCHF with no muon interactions, is where E4\_entry2 sits; run 3 passed through it in generation 1 and left.}
    \end{column}
  \end{columns}
\end{frame}




\begin{frame}{The new evolution found the next hole}
  \begin{columns}[c]
    \begin{column}{0.52\textwidth}\footnotesize
      Codex (OpenAI), 10 generations $\times$ 8 children on the full objective and the enriched dataset, 80 calls, 15 minutes. Its best: \key{reward $-17$}, against $-49$ for E4\_entry2.\\[1.5mm]
      \begin{tabular}{@{}lrr@{}}
        \toprule
        & E4\_entry2, 17.5\,mrad & evolved best\\
        \midrule
        planes & 5 (+3 second stations) & 15 \\
        cone & 17.5\,mrad & \key{50\,mrad}\\
        cost & 0.61\,MCHF & 7.2\,MCHF\\
        signal lost & 2.8\,\% & \key{74\,\%}\\
        $\mu$ / $\nu$ interactions & 0 / 0.62 & 0 / 0.06\\
        \bottomrule
      \end{tabular}
    \end{column}
    \begin{column}{0.45\textwidth}\footnotesize
      The signal tier saturates at 10\,\% loss (10 points); the neutrino floor costs $\approx$40 points through $100(1 - e^{-N})$. A huge cone vetoes neutrinos with \key{random hits}, exactly as it vetoes the signal: killing three quarters of the signal pays.\\[1.5mm]
      \takeaway[\footnotesize]{Again the optimiser is an adversary of the objective. The objective needs the signal loss unclipped (or as a constraint), or a sensitivity-like figure of merit (Punzi) where random-hit vetoes of background and signal cancel.}
    \end{column}
  \end{columns}
\end{frame}

\begin{frame}{An earlier hole (old run 1)}
  \begin{columns}[c]
    \begin{column}{0.52\textwidth}\footnotesize
      Run 1 used the \key{factorised} background only. Its best designs:\\[1.5mm]
      \begin{tabular}{@{}lrr@{}}
        \toprule
        E1 & factorised & direct\\
        \midrule
        false vertices / 15\,y & $5\cdot10^{-4}$ & $\mathbf{4.5\cdot10^{-3}}$\\
        surviving pairs & -- & 17 of 196\,736\\
        reward & $-0.50$ & $-0.89$\\
        \bottomrule
      \end{tabular}\\[2mm]
      Factorising assumes the two muons of a pair escape \key{independently}. With small fins they escape \key{together}: both enter through the same uncovered region. The optimiser trimmed into exactly that blind spot.
    \end{column}
    \begin{column}{0.45\textwidth}\footnotesize
      Fix: the larger of the two estimates. Run 2 then pays for the 17 pairs and finds a design that lets only 5 through (1 in the talk's sample), at lower cost.\\[1.5mm]
      Still in the model, not in the detector:
      \begin{itemize}\scriptsize
        \item no bounced muon of a pair is vetoed by E2; they are 0.7\,\% of the pairs
        \item $N$ rests on a handful of pairs: check on the bigger dataset
        \item fibre and pixel have no rate limit in \texttt{technology.py} (none in the documents)
        \item helium and air scattering, tilted fins: not yet
      \end{itemize}
      \takeaway[\footnotesize]{An optimiser is an adversary of the objective: every shortcut in the estimate becomes a design.}
    \end{column}
  \end{columns}
\end{frame}


% ======================================================================== PART 9
\blurb{What the simulation and the reconstruction still ignore, and what comes next.}
\section{Next}

\begin{frame}{What is not in yet}
  \begin{columns}[T]
    \begin{column}{0.48\textwidth}\small
      \key{Objective and optimisation}
      \begin{itemize}\footnotesize
        \item the signal tier saturates at 10\,\% loss: unclip it, or use a sensitivity (Punzi-like) figure of merit
        \item second stations as a design choice (layers per plane, with their cost and random hits)
        \item a new evolution on the fixed objective (Codex on the cluster, ready)
      \end{itemize}
      \key{Reconstruction}
      \begin{itemize}\footnotesize
        \item a dedicated search for the parent's entry hit, instead of a wide cone
        \item algorithm B: hit-driven tracklets (sees bounces); dead time and pile-up; tilted fins
      \end{itemize}
    \end{column}
    \begin{column}{0.48\textwidth}\small
      \key{Physics and data}
      \begin{itemize}\footnotesize
        \item statistics: the surviving muon interactions are a handful of simulated events (weight 0.5 each); the first metre beside the walls and the floor have $\sim$10 independent input muons per cell
        \item the rest of the enriched production (484 of 2\,000 seeds) and of our 20\,000-seed production
        \item photon conversions and secondary interactions (fast chain $-5$ to $-9$\,\% vs Geant4); shield and rock interactions
        \item specific signal models beyond the two-body decay
      \end{itemize}
    \end{column}
  \end{columns}
\end{frame}

\begin{frame}{Summary}
  \begin{columns}[T]
    \begin{column}{0.48\textwidth}\small
      \key{Data}: our own production keeps every muon category; an input enriched with a Pythia\,6 tail gives the high-energy SST1 muons $300\times$ more independent statistics.\\[2mm]
      \key{Backgrounds}: combinatorial pairs, muon interactions (Geant4 + Pythia\,6) and neutrino interactions (GENIE), all through one selection (pointing, two tracks) and one veto with a vertex cone.
    \end{column}
    \begin{column}{0.48\textwidth}\small
      \key{End to end}: the combinatorial background is gone ($<10^{-4}$); muon interactions in the walls and planes are always vetoed; in the helium they need the cone \emph{and} a second station at the entrance.\\[2mm]
      \takeaway[\footnotesize]{E4\_entry2 at 17.5\,mrad: 0.6 background events in 15 years, the neutrino floor; 2.8\,\% signal lost; 0.6\,MCHF. The optimiser's next run needs an objective that cannot trade signal for neutrinos.}
    \end{column}
  \end{columns}
\end{frame}

\end{document}
